

\hspace{-5em}\begin{minipage}{1.1\textwidth}
\begin{flushright}
{\fontsize{14}{22}\selectfont
\textbf{%
\dunderline[-8pt]{3pt}{PREFACE BY ROBERT GEFFNER}}
}
\end{flushright}
\end{minipage}


\vspace{6em}


\noindent{}In the past two decades, we have learned quite a bit about how family courts deal with allegations of intimate partner violence or abuse (IPV) and child maltreatment
(CM) in several countries, including the United
States. Since we have so much more research about the high prevalence
of IPV and the serious traumatic effects it produces in both the victims
and their children, it is surprising how frequently these issues are not
handled appropriately in family law child custody cases. Instead of following the guidelines of most professional organizations, such as the
National Council of Juvenile and Family Court Judges (NCJFC), 
Battered Women Justice Project (WJP), and the World Health Organization
(WHO), which are based upon careful and detailed research, we have all
too often seen that the IPV or CM reports and evidence are minimized,
discounted, or completely ignored in family court cases. If that was not
bad enough, too many mental health custody or parenting responsibility
evaluators recommend, and judges then actually award, custody to the
person that the victim or child has accused of abusing them. We have
seen these situations occur where there was a history of IPV by the domestic partner, the existence of prior restraining orders, and even where
prior partners of the offender were also abused in the same manner by
the offender. \pfl{}

We now also have a much better understanding of coercive control
and the abuse of power in these relationships, and how it is a major aspect of IPV. More jurisdictions are actually adding this factor into legal
statutes when a family court is dealing with issues of IPV. This also is
 consistent with the research that has been compiled over the past two
decades. In fact, it is the coercive control, intimidation, manipulation,
psychological abuse, and gaslighting that produces the most long-lasting
emotional damage and trauma to the victims of IPV and their children. \pfl{}

However, despite over two decades of research and practice concerning the negative and serious traumatic effects on these children who
grow up in homes with IPV, they still are not often believed when they
describe what they have experienced and observed, and why they do not
want to visit or have contact with the abusive parent, most often the father. The protective parent in these cases (usually the mother) is often
accused of \qq{alienating} the children against the other parent to counter
such allegations and disclosures made by the children. \pfl{}

Indeed, children have sometimes provided detailed disclosures of 
their having experienced 
physical, sexual, or emotional abuse to competent forensic 
interviewers. However, they are too often told
by family court professionals, or even child protective services at times,
that they are \qq{making it up} or that their parent (usually the mother but
sometimes the father) has coached or programmed them to say such
things in these cases. In such cases, research has shown that these accused fathers actually gain custody of a large number of the children
who had repeatedly stated that they do not want to live with, or sometimes not have contact with, the abusive parent. Labels are often used by
mental health custody evaluators, attorneys, guardians ad litem (GALs),
and judges to describe this situation where children do not want to have
contact with a parent they have accused of abuse or poor parenting. Such
terms or phrases as Resist and Refuse Dynamics (R-RD), gate keeping (GK), 
parent-child contact problems (PCCP), parental alienation
(PA), and so on, are often used by the family court professionals in these
situations. Labels and acronyms convert attitudes, beliefs, and behaviors
into \qq{things} or \qq{conditions} that become reified, subsequently taking
on a life of their own, while failing to account for a multitude of reasons
a child may not want to have contact with a parent, including abusive
behavior, an absent parent, or very poor parenting. \pfl{}\newpage{}

Many of these labels imply, assume, or actually suggest that the parent
who alleges IPV or reports the disclosures that their children made have
somehow programmed them to say these things. We have learned quite
a bit about suggestibility and coaching in the past two to three decades
from many methodologically sound studies, and the research supports
that it is quite rare for a child to be coached to make up detailed descriptions of sexually inappropriate behaviors in a child custody dispute. This
is so because the descriptions and disclosures by the children usually are
age and developmentally appropriate, remembered over time, and contain the details of abuse that a parent would not be able to replicate or
coach. Such details, expressions of strong emotions, and descriptions
of the sensations experienced by the abusive behaviors 
are eventually disclosed with good interviewing and therapeutic interactions. In fact,
some suggestibility research indicates that if a parent attempts to turn a
child against the other parent with whom they have had a good loving
relationship, it actually \qq{boomerangs} and the child resents the parent
trying to coach them to repeat negative statements about the other
parent. \pfl{}

\enlargethispage{1em}
Of course, the exception is when an abused child is sequestered in the
custody of the abuser for many years, and is totally cut off from 
the protective parent and any supportive people. Then, it is more common for
the child to \qq{accommodate} to the abuser because they have no one else
to support them. This accommodation is tragically seen in 
many protective mother cases when abused children are forced into the custody of
the abuser notwithstanding the child's vehement protestations 
to be protected from the abusive parent. The abused child feels abandoned and
turns against the protective parent after years of suffering both abuse
and abuse-denial by the abusive parent, and furthermore is angry at the
protective parent for not being able to really keep them safe from the
abusive parent. The abusive parent is aided in the denial of their abusive
behavior by the judge and mental health professionals who perpetuate
the myth that the abusive parent is indeed innocent even when there is
much evidence to the contrary. \pfl{}\newpage{}

It is interesting that the assumptions involved in such coaching scenarios imply that a parent untrained in psychological manipulation
techniques would be able to easily do such \qq{programming} of 
their children. If this was actually the case (which it is not), then almost all parents
would be able to raise children perfectly by just programming them to
do their chores, their homework, make their beds, clean their rooms,
and so on. Sometimes logic is completely ignored in these cases: biases
take hold, filter the information received, and influence perspectives
about the family, often to the detriment of the children. Any piece of
evidence or information is now fit into the biased perspectives, and thus
an uninformed narrative is now created. Once this bias 
leads to the psychological labels
previously noted, and decisions are made that are not in the best
interests of the children, then \qq{reunification} programs or \qq{camps} 
can be coercively forced upon these children 
to basically torture and scare them
to recant any allegations of abuse and to blame 
their mother presumably for making them say such things. Protective parents (who often
were the primary caretakers for years) may not see their children for
many years and not until they admit to coaching their children to make
up lies about the abusive parent. These programs have been debunked
as unethical and detrimental to children, and now some states in the
United States have statutes that prohibit courts from forcing children
into these programs. \pfl{}

\vspace*{0pt}\enlargethispage{3pt}\raggedbottom
So, the question becomes, with the expanding research and increase in knowledge over the past few decades concerning IPV and CM, how can
family courts, attorneys, and mental health evaluators continue to ignore or minimize IPV and CM in so many cases? There appears to
be two answers: The first is that people, including the professionals
just listed, do not want to admit or recognize the high prevalence of IPV and CM in families, even after so much public awareness
with stories saturating the media for many years now, and legislative bodies acting in response to these problems during the past two
decades. We now know about the millions of children sexually abused
by clergy, teachers, coaches, and parents in the last forty or more years.
However, believing a specific child who makes disclosures of sexual
abuse, or a woman who makes statements about experiencing IPV by
a person who appears to be respectable, \qq{upstanding,} and even 
a professional, does not fit into the expectations that many people have.
Societies have not yet accepted that these abusive events can occur in
families of all incomes, ethnicities, and status, even though they do
accept \qq{stranger} abuse, which is much less likely than what we are
discussing. \pfl{}

The second reason is the use of pseudo or \qq{junk} science that allows
people to have an excuse to both disbelieve the parental abuse during
a divorce, and to blame the victim instead. This started with the notion
of \qq{Parental Alienation Syndrome} (PAS) in the 1980s 
that was not research supported but, instead, was self-published and publicized as if
it were a real thing that was causing parents, mostly mothers, to make
up allegations of abuse to turn their children against the other parent,
mostly fathers, in divorce and child custody cases. This sounded good
for people who did not want to believe the high prevalence rates of CM
or IPV. As PAS was being discredited for numerous reasons in the late
1990s and early 2000s, wherein professional organizations and the official
diagnostic publications for mental and behavioral disorders were not accepting such untenable and baseless theories, it morphed into \qq{Parental
Alienation Disorder} (PAD).
However, research validating such a \qq{disorder} was 
suspect from the start, namely because its prevalence and
its specific diagnostic criteria were mostly unreliable, nonexistent, and
not validated. So, this permutation of PAS, just like its original form,
was again not accepted in any of the official professional diagnostic
categories nor by WHO. After more years had passed, PAS was morphed again by its proponents, who were often very well organized. As
such, it became just \qq{Parental Alienation} (PA), but with mostly the same
explanations and ideas as the original PAS. \pfl{}

However, PA, or just Alienation (A), are lay terms that people understand to mean a parent turning a child against the other parent (i.e.,
a negative behavior by a parent). Attempts at this can indeed occur
in dysfunctional families, so without paying attention to the pseudoscience of the way PAS/PAD/PA/A was being used, family courts and
the various professionals associated with them continued to use this
unsupported notion to minimize or counter allegations of abuse made
primarily by women who supported the reports made by their children.
It should be noted that this has also occurred with some fathers who
made good-faith allegations of child abuse committed by the mother, or
by the mother's boyfriend, that were countered by this pseudoscience. \pfl{}

It is a shame that, in this day and age, with all of the advancements
in technology and research, we are still so far behind in dealing with
IPV and CM in family courts and in society. Too many victims\mdash{}adults
and children\mdash{}have been revictimized by the use of pseudoscience and
the continued disbelief by too many family courts 
and associated professionals that IPV and/or CM can be happening at the levels that research
has shown for decades. It is much easier to believe that children and
adults are making things up, and to simultaneously blame the victim,
than to acknowledge the victimization and trauma experienced by so
many in families that break up. \pfl{}

In this second edition of \textit{From Madness to Mutiny}, Amy Neustein
and Michael Lesher follow up on their first book, highlighting many
of these issues and research findings that have come to the fore in the
past twenty years since their first edition was published. They provide
a detailed analysis of the family court system and how it has increasingly been influenced by cottage industries promoting pseudoscience
that have produced dire consequences for parents (mostly mothers) and
their children. The breadth of case histories and detailed longitudinal
case studies in this book bring into sharp focus the sad and recurring
problems in family courts in the United States and in other westernized
countries. This excellent book ties together the controversies, issues, and
the use of pseudoscience, and shows what has happened to far too many
women and children in family court cases when there are verifiable allegations of abusive behaviors. However, this book does not leave the
reader feeling that the situation is hopeless because the authors report
on ways the family court system can be changed to alleviate some of
its errors and reduce revictimization of the targets of IPV and CM. It is
important not only to identify the issues and reasons for the problems,
but to also make suggested changes that can lead to positive trauma-informed outcomes in these complex cases so that we can hopefully
finally see positive system changes. This book is an important effort
to illuminate the serious issues that have hindered family court decisions for years, and it provides a sensible roadmap for policy changes,
legislative actions, and social reform. \pfl{}

\hfill{}Robert Geffner, PhD, ABPP, ABN

\hfill{}Founding President, Family Violence \& Sexual

\hfill{}Assault Institute (FVSAI), dba Institute on Violence,

\hfill{}Abuse \& Trauma (IVAT) Distinguished Research Professor of

\hfill{}Psychology (Retired), Alliant International University

\hfill{}San Diego, CA

\hfill{}bgeffner@alliant.edu



